\documentclass[../HDF5_RFC.tex]{subfiles}

\begin{document}

%% ======================================================================
%%  Section 10 – Compatibility
%% ======================================================================
\section{Compatibility}
\label{sec:compatibility}

\subsection{Semantic Change in \texttt{H5Z\_class3\_t.name}}
\label{sec:name-semantic-change}

\noindent\colorbox{yellow!25}{\parbox{\dimexpr\linewidth-2\fboxsep}{%
  \textbf{Compatibility Note:} Upgrading a plugin from
  \texttt{H5Z\_class2\_t} to \texttt{H5Z\_class3\_t} changes the role of the
  \texttt{name} field from a freeform debug comment to a canonical identifier.
  The new \texttt{description} field in \texttt{H5Z\_class3\_t} preserves the
  display role: \texttt{H5Pget\_filter2} returns the \texttt{description} field
  for v3 plugins when it is non-\texttt{NULL}, ensuring backward-compatible
  output for applications that read the \texttt{name[]} output parameter.
  Plugin authors \textbf{should} populate \texttt{description} when upgrading
  to maintain continuity for downstream consumers.
}}\medskip

Upgrading a filter plugin from \texttt{H5Z\_class2\_t} to \texttt{H5Z\_class3\_t}
introduces no source-level incompatibility at existing \emph{call sites}
(code that calls \texttt{H5Zregister} continues to compile and link without
modification). However, the plugin binary itself \emph{is} binary incompatible:
the plugin shared library must be recompiled against the new struct definition,
and any existing binary that exports an \texttt{H5Z\_class2\_t} cannot simply be
reused as an \texttt{H5Z\_class3\_t} plugin without recompilation. Plugin
distributors must ship updated binaries. In addition, the upgrade entails a
\emph{semantic} shift in the role of the \texttt{name} field that plugin authors
and application maintainers must understand.

\verysubsection{Before (v2): freeform debug comment}
In \texttt{H5Z\_class2\_t}, the \texttt{name} field is documented as a ``Comment for
debugging'' and is also user-visible through the \texttt{name[]} output parameter of
\texttt{H5Pget\_filter2}. Existing plugins commonly populate this field with
human-readable prose such as \texttt{"HDF5 ZFP filter, version 1.0.0 (rate,
precision, and accuracy modes)"} or \texttt{"My Compression Library v2.3"}.

\verysubsection{After (v3): canonical identifier + description}
In \texttt{H5Z\_class3\_t}, the \texttt{name} field serves as the canonical identifier
used for registry lookup, on-demand plugin discovery, and as the string
accepted by \texttt{H5Pset\_filter\_by\_name}. It MUST conform to the naming rules in
\SectionRef{sec:param-rules}: printable ASCII, no commas, semicolons, or equals
signs, and ideally short and lowercase (e.g., \texttt{"zfp"}, \texttt{"blosc2"},
\texttt{"bitshuffle"}). Descriptive prose is not appropriate and MUST NOT be used.

The new \texttt{description} field in \texttt{H5Z\_class3\_t} preserves the display role
that the \texttt{name} field served in v2. It is an optional, human-readable string
(e.g., \texttt{"HDF5 ZFP filter, version 1.0.0 (rate, precision, and accuracy
modes)"}). No naming rules apply beyond requiring printable text. It is not
used for registry lookup or any programmatic dispatch.

\verysubsection{Impact on \texttt{H5Pget\_filter2}}
For v3 plugins, \texttt{H5Pget\_filter2} returns the \texttt{description} field in
its \texttt{name[]} output parameter when \texttt{description} is non-\texttt{NULL}.
If \texttt{description} is \texttt{NULL}, it falls back to the \texttt{name} field (the
canonical identifier). This means that plugin authors who populate
\texttt{description} during the v2-to-v3 upgrade preserve full backward
compatibility for applications that display the \texttt{H5Pget\_filter2} output.
Plugins that leave \texttt{description} as \texttt{NULL} will cause a
user-visible change in the \texttt{H5Pget\_filter2} output (the shorter canonical
name appearing where a longer description was previously shown). Depending on
what callers expect from \texttt{H5Pget\_filter2}, this may constitute an API
behavioral change even if the function signature is unchanged. Applications
that compare or display the \texttt{H5Pget\_filter2} string should be audited as
part of any v2-to-v3 plugin migration.

\verysubsection{Migration guidance}

\begin{itemize}
\item \textbf{Registry lookup and programmatic dispatch:} Use
  \texttt{H5Z\_class3\_t.name} (or the canonical name returned by
  \texttt{H5Pget\_filter\_name\_by\_idx}) as the primary identifier for all
  programmatic filter lookup. The canonical name is stable and suitable for
  storage, comparison, and round-tripping.
\item \textbf{Human-readable display:} Move the existing freeform v2
  \texttt{name} text into the \texttt{description} field when upgrading to v3.
  For example, a v2 plugin with
  \texttt{name = "HDF5 ZFP filter, version 1.0.0 (rate mode)"} should become
  a v3 plugin with \texttt{name = "zfp"} and
  \texttt{description = "HDF5 ZFP filter, version 1.0.0 (rate mode)"}.
  This ensures that \texttt{H5Pget\_filter2} continues to return the descriptive
  text. Filters that need to expose additional metadata (e.g., version strings)
  MAY also include those details as keys in the output of
  \texttt{H5Z\_get\_config\_func\_t}.
\item \textbf{No action required for consumers of \texttt{H5Pget\_filter2} that
  display the string verbatim:} These applications continue to function
  correctly, provided the plugin author populates the \texttt{description} field.
  Only applications that parse the freeform string to extract version metadata
  or descriptive information may need to adjust their display logic if the
  plugin author's choice of \texttt{description} text differs from the original
  v2 \texttt{name} text.
\end{itemize}

\subsection{On-Disk Format}

No changes. The string configuration is resolved entirely at configuration
time into \texttt{cd\_values}. The object header continues to store filter ID, flags,
and \texttt{cd\_values} as before. Files written using the string API are
byte-identical to files written using the existing \texttt{H5Pset\_filter} API with
equivalent parameters.

\subsection{Source Compatibility}

All existing public APIs are unchanged and not deprecated. Application code
using \texttt{H5Pset\_filter} continues to compile and function without modification.

\subsection{Binary / ABI Compatibility}

\texttt{H5Z\_class3\_t} is a new type, not a modification of \texttt{H5Z\_class2\_t}. New
functions are added to the public API. This is an API extension, not a
breaking change. A shared library minor version bump is appropriate.

\subsection{Plugin Compatibility}

Existing \texttt{H5Z\_class2\_t} plugins load and function without modification. They
can be used with the new API only via numeric ID with a raw \texttt{cd\_values}
parameter, or if they require no parameters. Plugin authors can upgrade to
\texttt{H5Z\_class3\_t} at their discretion.


%% ======================================================================
%%  Section 11 – Open Questions
%% ======================================================================
\section{Open Questions}
\label{sec:open-questions}

The following items require discussion and resolution before implementation
begins.

\verysubsection{Q1. Pipeline-string API}
Should a convenience function for specifying
entire pipelines in a single string be included in this release, or deferred?
See \SectionRef{sec:pipeline-string} for the tradeoffs. Reviewer feedback
strongly favors deferral, and this document recommends the same.

\verysubsection{Q2. Filter names in the HDF Group registry
  \textnormal{[Resolved: MUST be addressed before 1.0 release]}}
Canonical string names MUST be added to the HDF Group's external filter
registry as a normative field before the first release that includes
this feature. Without an authoritative name registry, independently
developed plugins may claim the same canonical name (e.g., two plugins both
registering \texttt{"lz4"}), and the first-registered-wins collision policy
provides no recourse after the fact. The following actions are required:
\begin{itemize}
\item The HDF Group MUST publish canonical names for all filters currently
  listed in the integer-ID filter registry~\cite{hdf5registeredfilters}
  alongside their numeric IDs.
\item The registration process for new filter IDs MUST be extended to
  require a canonical name that conforms to the naming rules in
  \SectionRef{sec:param-rules}.
\item The published canonical name is the single source of truth. Plugin
  authors MUST use the published name as their \texttt{H5Z\_class3\_t.name}
  field. Names that are not registered with the HDF Group are considered
  provisional and are subject to collision with future registrations.
\end{itemize}
This is a process requirement, not a technical change to the library.
Implementation may proceed in parallel, but the feature MUST NOT ship
without the registry update.

\verysubsection{Q3. API expansion}
The implementation introduces 7 new public or developer-facing functions and 3
new public types.  Each candidate for consolidation was considered; the
conclusions are summarized below.

\textbf{H5Pget\_filter\_name\_by\_idx + H5Pget\_filter\_params\_by\_idx.}
Merging these into a nine-parameter function was tried and reverted: the merged signature was unwieldy and callers who
only need the name (e.g., to dispatch on filter type) were forced to allocate a
params buffer they would never read.  The two-function design is retained.
Callers who genuinely need both MAY call them back-to-back; the index
validation overhead is negligible for pipeline sizes bounded by
\texttt{H5Z\_MAX\_NFILTERS} (32).

\textbf{H5Zlookup\_by\_name + H5Zfilter\_avail\_by\_name.}
These are intentionally asymmetric: \texttt{H5Zlookup\_by\_name} is a pure
in-memory query with no I/O side effects, while \texttt{H5Zfilter\_avail\_by\_name}
may trigger a filesystem scan and shared-library \texttt{dlopen}.  Collapsing them
into a single function with a \texttt{trigger\_load} flag would hide a potentially
expensive, irreversible side effect behind a boolean parameter---a pattern the
HDF5 API avoids elsewhere for good reason.  The two functions are retained as
separate entry points with clearly distinct semantics.

\textbf{H5Zregister\_name + H5Zpreregister\_name.}
These serve orthogonal use cases: \texttt{H5Zregister\_name} binds an alias
immediately (requires the filter to be present), while
\texttt{H5Zpreregister\_name} supports deferred binding for initialization
sequences that run before plugin load.  A single function could not satisfy
both contracts without an ordering dependency that would be invisible at the
call site.  They are retained as separate functions.

\textbf{H5Zconfig\_get\_param.}
This function is placed in \texttt{H5Zdevelop.h} (developer-facing) rather than
the public API, limiting the public surface.  Plugin authors who do not use
it incur no dependency.

\textbf{Net assessment.}  The 7-function surface is the result of deliberate
prior consolidation (R1) and clear semantic separation.  No further reductions
are identified that would not sacrifice correctness, performance predictability,
or API clarity.  If future releases warrant it, \texttt{H5Pget\_filter\_name\_by\_idx}
and \texttt{H5Pget\_filter\_params\_by\_idx} could be supplemented by a
convenience wrapper that fills both buffers in one call, but the two focused
functions MUST remain available as the normative API.

%% ======================================================================
%%  Section 12 – Test Plan
%% ======================================================================
\section{Test Plan}
\label{sec:test-plan}

The following test categories MUST be satisfied before the implementation is
merged.  Test names in \texttt{typewriter} font are illustrative identifiers
for the test matrix; the actual test suite may group them differently.

\subsection{Parser Unit Tests}
\label{sec:test-parser}

\noindent The parameter string parser (\SectionRef{sec:param-rules}) is the
lowest-level component and MUST be exercised independently of any filter
registration.  Tests invoke \texttt{H5Zconfig\_get\_param} (\SectionRef{sec:config-get-param})
directly with hand-crafted strings.

\verysubsection{Valid inputs}

\begin{center}
\small
\renewcommand{\arraystretch}{1.25}
\begin{tabularx}{\textwidth}{l>{\raggedright\arraybackslash\hsize=0.7\hsize}X>{\raggedright\arraybackslash\hsize=1.3\hsize}X}
\toprule
\textbf{Test ID} & \textbf{Input} & \textbf{Expected result} \\
\midrule
\texttt{parse-01} & \texttt{"level=6"}
  & key=\texttt{level}, value=\texttt{"6"} \\
\texttt{parse-02} & \texttt{"mode=rate, rate=3.5"}
  & two key-value pairs; trailing whitespace stripped \\
\texttt{parse-03} & \texttt{"fast\_mode, level=6"}
  & bare key \texttt{fast\_mode} (empty value) + \texttt{level=6} \\
\texttt{parse-04} & \texttt{"level="}
  & key=\texttt{level}, explicit empty string value \\
\texttt{parse-05} & \texttt{NULL}
  & valid; zero parameters \\
\texttt{parse-06} & \texttt{""}
  & valid; zero parameters \\
\texttt{parse-07} & \texttt{'dict\_path="/data/run\_1,v2/dict.bin"'}
  & comma inside quotes is literal; value=\texttt{/data/run\_1,v2/dict.bin} \\
\texttt{parse-08} & \texttt{'note="say ""hi"""'}
  & doubled double-quote escape; value=\texttt{say "hi"} \\
\texttt{parse-09} & \texttt{"  level  =  6  "}
  & whitespace stripped; key=\texttt{level}, value=\texttt{6} \\
\texttt{parse-10} & \texttt{'level = " value "'}
  & key=\texttt{level}; value=\texttt{\ value\ } (leading/trailing spaces preserved inside quoted value) \\
\texttt{parse-11} & \texttt{"LEVEL=6"}
  & key normalized to lowercase \texttt{level}; value preserved as \texttt{"6"} \\
\bottomrule
\end{tabularx}
\end{center}

\verysubsection{Malformed inputs (must return \texttt{H5E\_BADVALUE})}

\begin{center}
\small
\renewcommand{\arraystretch}{1.25}
\begin{tabularx}{\textwidth}{lX}
\toprule
\textbf{Test ID} & \textbf{Input and reason for rejection} \\
\midrule
\texttt{parse-E01} & \texttt{"=value"} --- empty key before \texttt{=} \\
\texttt{parse-E02} & \texttt{",level=6"} --- leading comma produces empty first token \\
\texttt{parse-E03} & \texttt{"level=6,"} --- trailing comma produces empty last token \\
\texttt{parse-E04} & \texttt{'path="/no/closing/quote'} --- unbalanced opening double-quote \\
\texttt{parse-E05} & \texttt{"level=6, level=6"} --- duplicate key \\
\texttt{parse-E06} & String longer than \texttt{H5Z\_CONFIG\_STRING\_MAX} (4096) bytes --- exceeds maximum \\
\texttt{parse-E07} & String with more than \texttt{H5Z\_CONFIG\_MAX\_PARAMS} (64) comma-separated tokens --- exceeds maximum parameter count \\
\bottomrule
\end{tabularx}
\end{center}

\verysubsection{Boundary conditions}

\begin{itemize}
\item \texttt{parse-B01}: Single character key and value (\texttt{"a=1"}).
\item \texttt{parse-B02}: Maximum allowed string length
  (\texttt{H5Z\_CONFIG\_STRING\_MAX} = 4096 bytes) --- accepted.
\item \texttt{parse-B03}: \texttt{H5Z\_CONFIG\_STRING\_MAX}+1 byte string --- rejected
  with \texttt{H5E\_BADVALUE}.
\item \texttt{parse-B04}: Key consisting entirely of printable non-special ASCII
  (\texttt{"\_0123456789abcdefghijklmnopqrstuvwxyz=ok"}).
\item \texttt{parse-B05}: Value containing semicolons inside double quotes --- accepted;
  value containing a bare semicolon outside quotes --- rejected (reserved delimiter).
\item \texttt{parse-B06}: Exactly \texttt{H5Z\_CONFIG\_MAX\_PARAMS} (64) parameters ---
  accepted.
\item \texttt{parse-B07}: \texttt{H5Z\_CONFIG\_MAX\_PARAMS}+1 (65) parameters --- rejected
  with \texttt{H5E\_BADVALUE}.
\end{itemize}


\subsection{Callback Contract Tests}
\label{sec:test-callbacks}

These tests verify the two-pass \texttt{set\_config} contract
(\SectionRef{sec:callback-specs}) and the \texttt{get\_config} round-trip guarantee.

\begin{itemize}
\item \texttt{cb-01}: First-pass call with \texttt{cd\_values=NULL} returns the correct
  \texttt{cd\_nelmts}; second-pass call populates \texttt{cd\_values} without exceeding
  the allocated buffer.
\item \texttt{cb-02}: A mock callback that returns a larger \texttt{cd\_nelmts} on the
  second pass triggers a library error (not a buffer overrun).
\item \texttt{cb-03}: A mock callback that returns a smaller \texttt{cd\_nelmts} on the
  second pass triggers a library error (not silent truncation).
\item \texttt{cb-04}: A callback that modifies \texttt{flags} does so silently;
  verify the updated flags value by comparing the input flags with those
  returned by \texttt{H5Pget\_filter2} after the call.
\item \texttt{cb-05}: A callback that does NOT modify \texttt{flags} leaves the
  error stack empty after the call.
\item \texttt{cb-06}: \texttt{get\_config} called with \texttt{buf=NULL} returns the required
  length without writing any bytes.
\item \texttt{cb-07}: \texttt{get\_config} output is a valid \texttt{set\_config} input (round-trip
  grammar check: parse the output and confirm it tokenizes without error).
\end{itemize}


\subsection{Built-in Filter Round-Trip Tests}
\label{sec:test-builtin-roundtrip}

For each built-in filter that implements \texttt{set\_config}/\texttt{get\_config}, a
round-trip test confirms that:
\begin{align*}
  \text{set\_config}(s) \;\Rightarrow\; cd\_values_1
  \quad\text{and}\quad
  \text{get\_config}(cd\_values_1) \;\Rightarrow\; s'
  \quad\text{and}\quad
  \text{set\_config}(s') \;\Rightarrow\; cd\_values_2
  \quad\text{with}\quad cd\_values_1 = cd\_values_2
\end{align*}
The reconstructed string $s'$ need not equal $s$ literally; only the
resulting \texttt{cd\_values} must be identical.

\begin{center}
\small
\renewcommand{\arraystretch}{1.25}
\begin{tabular}{lll}
\toprule
\textbf{Test ID} & \textbf{Filter} & \textbf{Representative parameter string} \\
\midrule
\texttt{rt-01} & deflate    & \texttt{"level=0"} (minimum) \\
\texttt{rt-02} & deflate    & \texttt{"level=6"} (default) \\
\texttt{rt-03} & deflate    & \texttt{"level=9"} (maximum) \\
\texttt{rt-04} & szip       & \texttt{"coding=entropy, pixels\_per\_block=2"} \\
\texttt{rt-05} & szip       & \texttt{"coding=nn, pixels\_per\_block=32"} \\
\texttt{rt-06} & scaleoffset & \texttt{"scale\_type=float\_dscale, scale\_factor=4"} \\
\texttt{rt-07} & scaleoffset & \texttt{"scale\_type=float\_escale, scale\_factor=0"} \\
\texttt{rt-08} & scaleoffset & \texttt{"scale\_type=int, scale\_factor=8"} \\
\texttt{rt-09} & shuffle    & \texttt{NULL} (parameterless) \\
\texttt{rt-10} & fletcher32 & \texttt{NULL} (parameterless) \\
\texttt{rt-11} & nbit       & \texttt{NULL} (parameterless) \\
\bottomrule
\end{tabular}
\end{center}

\noindent Additional round-trip tests MUST verify that:
\begin{itemize}
\item \texttt{rt-E01}: Passing a non-NULL, non-empty \texttt{params} to a parameterless
  built-in (\texttt{shuffle}, \texttt{fletcher32}, \texttt{nbit}) returns \texttt{H5E\_BADVALUE}.
\item \texttt{rt-E02}: Passing an unrecognized key (e.g., \texttt{"clevl=6"}
  instead of \texttt{"clevel=6"}) to any built-in returns
  \texttt{H5E\_BADVALUE} with an error message identifying the unrecognized key.
  v3 filters MUST reject unknown keys (\SectionRef{sec:callback-specs}).
\item \texttt{rt-E03}: Passing an out-of-range value (\texttt{"level=10"} for deflate,
  \texttt{"pixels\_per\_block=33"} for szip) returns \texttt{H5E\_BADVALUE}.
\end{itemize}


\subsection{Regression Tests: v1/v2 Plugin Behavior Unchanged}
\label{sec:test-regression}

These tests confirm that no existing behavior is broken.

\begin{itemize}
\item \texttt{reg-01}: A filter registered as \texttt{H5Z\_class2\_t} remains callable via
  \texttt{H5Pset\_filter} with the same integer ID and \texttt{cd\_values} as before.
  Data written and read through such a filter before and after the feature
  introduction MUST be byte-identical.
\item \texttt{reg-02}: \texttt{H5Pset\_filter\_by\_name} with \texttt{params=NULL} for a v2 filter
  succeeds (NULL is treated as ``no parameters'').
\item \texttt{reg-03}: \texttt{H5Pset\_filter\_by\_name} with a non-NULL \texttt{params} for a v2 filter
  that does NOT implement \texttt{set\_config} returns \texttt{H5E\_UNSUPPORTED}.
\item \texttt{reg-04a}: The existing \texttt{H5Pget\_filter2} continues to return the
  \texttt{name} field from the class struct (the debug comment) for v2 plugins,
  not the canonical registry name.  Verify that a v2 plugin with a verbose
  name field (e.g., \texttt{"HDF5 deflate filter v1.0"}) returns that string,
  not \texttt{"deflate"}.
\item \texttt{reg-04b}: For a v3 plugin with a non-\texttt{NULL} \texttt{description}
  field, \texttt{H5Pget\_filter2} returns the \texttt{description} text, not the
  canonical \texttt{name}. For example, a plugin with \texttt{name = "zfp"} and
  \texttt{description = "HDF5 ZFP filter v1.0"} returns
  \texttt{"HDF5 ZFP filter v1.0"}.
\item \texttt{reg-04c}: For a v3 plugin with \texttt{description = NULL},
  \texttt{H5Pget\_filter2} falls back to returning the canonical \texttt{name}
  field (e.g., \texttt{"zfp"}).
\item \texttt{reg-05}: All existing \texttt{h5repack} invocations with the legacy filter
  syntax (\texttt{GZIP}, \texttt{SHUF}, \texttt{FLET}, \texttt{UD=id,...}) produce
  output files byte-identical to those produced by the pre-feature version of
  \texttt{h5repack}.
\end{itemize}


\subsection{Parallel HDF5 Tests}
\label{sec:test-parallel}

These tests require a Parallel HDF5 build and at least two MPI ranks.

\begin{itemize}
\item \texttt{par-01} (\textbf{DCPL consistency}): All ranks call
  \texttt{H5Pset\_filter\_by\_name} with the same name and parameter string; verify
  that the resulting \texttt{cd\_values} in the DCPL are byte-identical across ranks
  before \texttt{H5Dcreate}.
\item \texttt{par-02} (\textbf{rank failure propagation}): Simulate a missing plugin
  on one rank (e.g., by unsetting \texttt{HDF5\_PLUGIN\_PATH} on rank~1 only).
  Verify that the failing rank returns \texttt{H5E\_NOFILTER} from
  \texttt{H5Pset\_filter\_by\_name} and that the application can detect this and
  abort collectively without an MPI deadlock.
\item \texttt{par-03} (\textbf{debug-build consistency check}): In a debug build
  (\texttt{NDEBUG} not defined), induce a rank-inconsistent DCPL (by having one rank
  add an extra filter entry) and verify that \texttt{H5Dcreate} returns
  \texttt{H5E\_CANTCREATE} with an informative diagnostic rather than silently
  producing a corrupt file.
\item \texttt{par-04} (\textbf{built-in-only pipeline}): Create a chunked dataset using
  only built-in filters configured via the string API in a multi-rank parallel
  job; verify that the dataset can be written and read back correctly from all
  ranks.
\end{itemize}


\subsection{Thread-Safety Tests}
\label{sec:test-threadsafety}

These tests require a thread-safe HDF5 build (\texttt{--enable-threadsafe}).

\begin{itemize}
\item \texttt{ts-01} (\textbf{concurrent lookups}): Spawn $N$ threads
  ($N \geq 8$) that each call \texttt{H5Zlookup\_by\_name} and
  \texttt{H5Zfilter\_avail\_by\_name} for built-in filter names in a tight loop.
  No data race or crash MUST occur.
\item \texttt{ts-02} (\textbf{concurrent registration and lookup}): One thread calls
  \texttt{H5Zregister\_name} for a v2 plugin while other threads concurrently call
  \texttt{H5Zlookup\_by\_name}.  The lookup MUST see either the old state (name not
  present) or the new state (name present); no partial or corrupt state.
\item \texttt{ts-03} (\textbf{concurrent \texttt{H5Pset\_filter\_by\_name}}): Threads each
  create their own DCPL and call \texttt{H5Pset\_filter\_by\_name} concurrently for the
  same built-in filter name.  All calls MUST succeed without error or data
  corruption (property lists are per-object and not shared).
\item \texttt{ts-04} (\textbf{on-demand plugin scan under contention}): Arrange for a
  plugin to be loaded on-demand (remove its entry from the registry first).
  Then have $N$ threads simultaneously call \texttt{H5Zfilter\_avail\_by\_name} for that
  name.  Verify the plugin is registered exactly once and all threads receive a
  positive result.
\end{itemize}


\subsection{Name Registry Tests}
\label{sec:test-registry}

\begin{itemize}
\item \texttt{nr-01}: \texttt{H5Zregister\_name} for a name that already exists returns
  an error (first-registered-wins).
\item \texttt{nr-02}: Multiple names may map to the same filter ID (aliases);
  \texttt{H5Zlookup\_by\_name} returns the same ID for all aliases.
\item \texttt{nr-03}: \texttt{H5Zlookup\_by\_name} lookup is case-insensitive
  (\texttt{"Deflate"}, \texttt{"DEFLATE"}, and \texttt{"deflate"} all return
  \texttt{H5Z\_FILTER\_DEFLATE}).
\item \texttt{nr-04}: After \texttt{H5Zunregister}, all name registry entries for the
  unregistered filter ID are removed; subsequent \texttt{H5Zlookup\_by\_name} returns
  \texttt{H5Z\_FILTER\_ERROR}.
\item \texttt{nr-05}: \texttt{H5Zpreregister\_name} followed by \texttt{H5Zregister} promotes
  the pending entry into the live registry.
\item \texttt{nr-06}: \texttt{H5Zpreregister\_name} entries for filters that are never
  registered are discarded silently at \texttt{H5close} (no memory leak; verify
  with a memory checker).
\item \texttt{nr-07}: \texttt{H5Zregister\_name} called before \texttt{H5Zregister} returns
  \texttt{H5E\_NOTFOUND}.
\item \texttt{nr-08}: A v3 plugin whose \texttt{name} field violates the naming rules
  (\SectionRef{sec:param-rules}) is rejected at \texttt{H5Zregister} time with
  \texttt{H5E\_BADVALUE}; the filter itself is NOT added to the filter table.
\end{itemize}


%% ======================================================================
%%  Section 13 – Risks
%% ======================================================================
\section{Risks}
\label{sec:risks}

\verysubsection{Grammar maintenance}
The key=value format with quoted-value support and
bare-key flags constitutes a specification the library must support
indefinitely. Future grammar extensions MUST be backward-compatible: any
string valid under the current grammar MUST remain valid with the same
semantics under future grammars. The semicolon is reserved in anticipation of
a possible pipeline-string extension.

\verysubsection{Plugin adoption}
The value of the string API scales with the number of
filters that implement \texttt{set\_config}. If adoption is slow, v2 plugins can only
be used via the existing \texttt{H5Pset\_filter} API. Clear documentation and a low
barrier to v3 adoption mitigate this.

\verysubsection{Name collisions}
Independently developed plugins could choose the same
canonical name. The first-registered-wins collision policy prevents silent
breakage at runtime. This risk is mitigated by the requirement that canonical
names be published in the HDF Group's central filter
registry~\cite{hdf5registeredfilters} before this feature ships (see Q2 in
\SectionRef{sec:open-questions}). Plugins that register names not listed in the
central registry are considered provisional and may collide with future
registrations.

\verysubsection{Plugin scan cost}
Name-based plugin discovery scans all shared libraries
in \texttt{HDF5\_PLUGIN\_PATH} until a match is found. This is not a new cost---it
is identical to how the existing integer-ID plugin search already works for
all plugin types (filters, VFDs, VOLs). The name registry caches the result
so the scan occurs only on first reference to an unknown name.

\verysubsection{Parsing overhead for small-data workloads}
The string API adds parsing cost at \texttt{H5Pset\_filter\_by\_name} call time:
one hash-table lookup for name resolution and one linear pass over the
parameter string (bounded by \texttt{H5Z\_CONFIG\_STRING\_MAX} = 4096~bytes) per
filter, followed by two invocations of the \texttt{set\_config} callback (the
two-pass size-query protocol). For typical scientific workloads where dataset
creation is infrequent relative to I/O, this overhead is negligible.

In ``small data'' workloads that create thousands of lightweight datasets in a
tight loop, the per-dataset overhead of string parsing could become measurable,
though the dominant cost is already property list manipulation and object header
creation.

Applications that are sensitive to dataset-creation latency and configure the
same filter pipeline repeatedly SHOULD construct the DCPL once using
\texttt{H5Pset\_filter\_by\_name} and reuse it across multiple \texttt{H5Dcreate}
calls, avoiding repeated parsing. This is already the recommended practice for
\texttt{H5Pset\_filter} and does not change with the new API.

\noindent\colorbox{yellow!25}{\parbox{\dimexpr\linewidth-2\fboxsep}{%
  \textbf{Open Design Decision --- Pre-tokenization on the DCPL:}
  An alternative to the current design---where each \texttt{set\_config} callback
  re-parses the raw parameter string via \texttt{H5Zconfig\_get\_param}---is to
  perform tokenization once at the point \texttt{H5Pset\_filter\_by\_name} is called
  and store the resulting key-value token array on the DCPL.  The
  \texttt{set\_config} callback would then receive pre-parsed tokens rather than
  the raw string, eliminating redundant parsing.  This also becomes necessary
  if the pipeline-string extension (\SectionRef{sec:pipeline-string}) is
  adopted: a single string encoding parameters for multiple filters requires
  the library to tokenize and dispatch tokens to the correct filter anyway,
  so the callback interface would need to change to accept tokens regardless.
  The group should decide whether to revise the callback interface now to
  accept pre-parsed tokens, or retain the current raw-string interface and
  revisit if the pipeline-string extension is adopted.
}}

\end{document}
